\scanpage{563}
\mspage{474}
\deleted{1°) Catégorie définie par un objet $x$ de $C$, un
projecteur $p$ de $x$ tel que $p\ne\operatorname{id}_x$, un
objet $x^{\unsure{\leftrightarrows|\natural}}$ et un couple \deleted{\unsure{de}} de flèches
\diagram{\begin{tikzcd}[ampersand replacement=\&]
x^{\leftrightarrows}\arrow[r,shift left,"\sigma"]\&x\arrow[l,shift left,"\rho"]
\end{tikzcd}}
faisant de $x^{\leftrightarrows}$ une image de $p$,
enfin une flèche $x\xrightarrow{u}y$ telle que
$up=u$, et la flèche composée}

[On prend pour $C'$ la catégorie
\diagram{\small\(\displaystyle \unsure{Q|\mathcal Q}=\left(\vcenter{\hbox{\begin{tikzcd}[ampersand replacement=\&,column sep=large,row sep=large]
{\unsure{\xi|\zeta}_0=\xi^{\leftrightarrows}}\arrow[r,shift left,"\sigma"]\arrow[dr,"u_0"']\&\xi\arrow[l,shift left,"\rho"]\arrow[loop right,"p"]\arrow[d,"u=u_0\rho"]\\
{}\& y
\end{tikzcd}}}\quad
\begin{array}{ll}
\rho\sigma=\operatorname{id}_{\deleted{\unsure{\xi|x}}\unsure{\xi_0}}&\\
\sigma\rho=p&\text{d'où }\left|\begin{aligned}p^2&=p\\p\sigma&=\sigma\\\rho p&=\rho\end{aligned}\right.\\
u=u_0\rho&\text{d'où }\left|\begin{aligned}up&=u\\u\sigma&=u_0\end{aligned}\right.\\
\deleted{\unsure{u\rho=}}&
\end{array}\right)\)}
(trois objets $\xi,\xi_0,y$ et cinq flèches non identiques
$\sigma,\rho,\unsure{p},u_0,u$, liées par les relations ci-dessus),
ou la catégorie (isomorphe à $Q^{\circ}$)
\diagram{\small\(\displaystyle Q'=\left(\vcenter{\hbox{\begin{tikzcd}[ampersand replacement=\&,column sep=large,row sep=large]
\xi_0\arrow[r,shift left,"\sigma"]\&\xi\arrow[l,shift left,"\rho"]\arrow[loop right,"p"]\\
{}\& y\arrow[u,"u"']\arrow[ul,"u_0"]
\end{tikzcd}}}\quad
\begin{array}{l}
\left.\begin{aligned}\rho\sigma&=\operatorname{id}_{\xi_0}\\\sigma\rho&=p\end{aligned}\right\}
\Rightarrow\left|\begin{aligned}p^2&=p\\p\sigma&=\sigma\\\rho p&=\rho\end{aligned}\right.\\
u=\sigma u_0\ \Rightarrow\left|\begin{aligned}pu&=u\\\rho u&=u_0\end{aligned}\right.
\end{array}\right)\)}
On \deleted{prend des}\inserted{se limite aux} foncteurs $Q\longrightarrow C$, $Q'\mathrel{\unsure{-|\longrightarrow}}C$
tels que$^{\unsure{*|!}}$ $p$ soit transformé en une
flèche \uline{non identique} de $C^{*}$ (donc la
% slanted
\margin{je ne suis pas sûr si c'est raisonnable --- à vérifier avec \unsure{soin} !}
condition relative à $Q,Q'$ devient vide, si
$C$ est sans projecteurs).]

Sous ces conditions, $\unsure{\mathcal E|E}$ est accessible, et
$f:\mathcal E\mathrel{-}C$ est accessible.
% archive sheet number: 562
